\section{Joint dimension bounds for operational descriptions}
\label{sec:interiorcoding79}
The dimension-free comparison of Lemma~\ref{lem:interiormetric78}
allows a description theorem whose constants are uniform in the
matrix dimension.  This complements, rather than replaces, the
closed-body law of Theorem~\ref{thm:matrixcover76}.  Put
\[
 \mathfrak I_d=\{E=E^*:I_d/4\preceq E\preceq3I_d/4\},\qquad p=d^2.
\]
Throughout this section, centres are legal ordered binary effects
and the operational distance is the unhalved distance $d_N$ of the
same consuming, classical-output interface.  In particular, the
future tester may retain a reference, use coherent adaptive controls,
and stop at a public bounded time.  All logarithms in description
lengths have base two.

\begin{theorem}[Uniform interior operational entropy]
\label{thm:interiorentropy79}
Let $d,N\ge1$ and $0<\delta\le2^{-13}$.  The least cardinality
$\mathcal C_N(\mathfrak I_d,\delta)$ of a $d_N$-cover of
$\mathfrak I_d$, allowing centres anywhere in $\mathfrak E_d$, satisfies
\begin{equation}\label{eq:interiorentropy79}
 \left(\frac{\sqrt N}{2048\delta}\right)^{p}
 \le \mathcal C_N(\mathfrak I_d,\delta)
 \le \left(\frac{25\sqrt N}{\delta}\right)^{p}.
\end{equation}
The upper bound is attained by a public finite exact rational
construction when $\delta$ is rational.  For real $\delta$, a
volumetric net gives the stated upper bound.  Consequently the
optimal deterministic fixed-length description into legal effects is
\begin{equation}\label{eq:interiorbits79}
 B_{\rm int}^{\star}(d,N,\delta)
 =\frac{d^2}{2}\log_2 N+d^2\log_2(1/\delta)+O(d^2),
\end{equation}
where the remainder is uniform in all three parameters.
\end{theorem}
\begin{proof}
The real vector space of Hermitian matrices has dimension $p$, and
$\mathfrak I_d$ is precisely the translate of its operator-norm ball
of radius $1/4$.  Set $s=512\delta/\sqrt N$.  A maximal subset
$\{E_1,\ldots,E_L\}$ of $\mathfrak I_d$ whose distinct points have
operator distance strictly greater than $s$ is finite and its closed
$s$-balls cover $\mathfrak I_d$.  Translation-invariant Lebesgue
measure and homogeneity therefore give
\[
 Ls^p\ge(1/4)^p,
 \qquad L\ge\left(\frac{\sqrt N}{2048\delta}\right)^p.
\]
The unit-ball volume cancels; no dimension-dependent volume estimate
is used.  Lemma~\ref{lem:interiormetric78} separates any two selected
effects in $d_N$ by at least $4\delta$, since $512\delta\le1/16$.
A closed operational ball of radius $\delta$ contains at most one
selected effect, irrespective of its centre.  This uses the triangle
inequality for $d_N$, not a property of a comparison modulus.

For the upper bound choose a maximal operator-norm
$\delta/(4\sqrt N)$-separated subset of $\mathfrak I_d$.
It covers the latter at that operator radius and hence at operational
radius $\delta$.  The disjoint open balls of half that radius lie
inside the concentric ball of radius
$1/4+\delta/(8\sqrt N)$.  Their number is at most
$(1+2\sqrt N/\delta)^p$, which is smaller than the asserted bound.
The rational construction with a specified finite stopping rule is
proved next.  A fixed public $B$-bit decoder has at most $2^B$
outputs, and indexing a smallest finite cover suffices.  Taking
logarithms, with one bit for rounding to an integer, proves
\eqref{eq:interiorbits79}.
\end{proof}

\subsection{An exact public interior dictionary}
Only the public parameters $(d,N,\delta)$ determine the dictionary.
Let $\delta\in\Q\cap(0,1]$ and set
\begin{equation}\label{eq:interiorgrid79}
 k=\lceil\sqrt N\rceil,\qquad
 K=\left\lceil\frac{128dk}{\delta}\right\rceil,\qquad
 t=\frac{\delta}{16k}.
\end{equation}
The integer $k$ is computed by an integer square-root test.  Enumerate
Hermitian matrices with diagonal coordinates in
$\{0,1/K,\ldots,1\}$ and real and imaginary off-diagonal coordinates
in $\{-1,-1+1/K,\ldots,1\}$.  Order coordinates as in
Theorem~\ref{thm:matrixcodec77} and retain only the matrices in
$[I_d/8,7I_d/8]$.  From this ordered finite stream construct a list
$\mathcal D$ greedily, accepting $G$ precisely when
$\|G-C\|_{\rm op}>t$ for every previously accepted $C$.
For a Hermitian rational difference $H$, the complementary test is
\begin{equation}\label{eq:interiortest79}
 \|H\|_{\rm op}\le t
 \quad\Longleftrightarrow\quad tI_d-H\succeq0
                      \ \hbox{and}\ tI_d+H\succeq0.
\end{equation}
Thus exact rational positive-semidefinite tests settle the comparison,
including equality, without computing eigenvalues.  The Schur test
of Theorem~\ref{thm:matrixcodec77} applies to both inequalities.

\begin{theorem}[Finite exact interior coding]
\label{thm:interiorcodec79}
For every $d,N\ge1$ and rational $0<\delta\le1$, the preceding
construction terminates and defines a common encoder and decoder.
For a supplied Gaussian-rational $E\in\mathfrak I_d$, the decoded
centre $C\in[I_d/8,7I_d/8]$ satisfies
\begin{equation}\label{eq:interiorcodeerror79}
 d_N(E,C)\le5\delta/16.
\end{equation}
The same conclusion holds for a real $E\in\mathfrak I_d$ supplied
with a Gaussian-rational Hermitian approximation $A$ certified to
satisfy $\|A-E\|_{\rm op}\le d/K$.  The payload is an index of
length $\lceil\log_2|\mathcal D|\rceil$, and
\begin{equation}\label{eq:interiorcodesize79}
 |\mathcal D|\le(1+12k/\delta)^p
                  \le(25\sqrt N/\delta)^p.
\end{equation}
All words of that length have legal decodings.  No nonpublic
matrix coordinates or target-dependent dictionary are free advice.
\end{theorem}
\begin{proof}
Round every real coordinate of $A$ to the nearest multiple of $1/K$,
with ties towards positive infinity, and preserve conjugate entries.
Call the result $G$.  In every row the diagonal rounding error is at
most $1/(2K)$ and each off-diagonal error has modulus at most
$1/(\sqrt2K)$.  The Hermitian row-sum bound gives
$\|G-A\|_{\rm op}\le d/K$.  Therefore
\[
 \|G-E\|_{\rm op}\le2d/K\le\delta/(64k)\le1/64.
\]
For rational $E$ take $A=E$; the same, slightly looser, bound applies.
Weyl's inequality puts $G$ in $[I_d/8,7I_d/8]$.  Its diagonal and
off-diagonal coordinates also lie in the enumerated ranges, so it
occurs in the legal finite stream.

If $G$ is retained, select $C=G$.  Otherwise the strict greedy rule
supplies an earlier retained centre with $\|G-C\|_{\rm op}\le t$;
select the first such centre.  Hence
\[
 \|E-C\|_{\rm op}\le\frac{\delta}{64k}
                            +\frac{\delta}{16k}
                      =\frac{5\delta}{64k}.
\]
Both endpoints lie in the range of Lemma~\ref{lem:interiormetric78},
which proves \eqref{eq:interiorcodeerror79}.  Legality and both
selection tests use finite exact rational arithmetic.

Distinct dictionary elements have operator distance greater than
$t$, and every element lies within operator radius $3/8$ of $I_d/2$.
The open operator balls of radius $t/2$ about them are disjoint and
lie in the ball of radius $3/8+t/2$.  Cancelling the common unit-ball
volume gives $(1+3/(4t))^p=(1+12k/\delta)^p$.  Since
$k\le2\sqrt N$ and $\sqrt N/\delta\ge1$, the second inequality
in \eqref{eq:interiorcodesize79} follows.

The decoder reconstructs the entire ordered list from the public
parameters before interpreting its index.  Unused fixed-length
words return $I_d/2$.  The list is nonempty because the rounded
centre $I_d/2$ occurs.  This proves termination and total legal
decoding.  A rational approximation to a real input is supplied
data under the displayed certificate, not an arbitrary exact-real
oracle provided by the theorem.
\end{proof}

The raw enumeration has at most
$(K+1)^d(2K+1)^{d(d-1)}$ tuples.  Exact comparisons against the
retained list and the bit lengths of Schur-complement entries are
additional reconstruction costs.  They are not controlled by
\eqref{eq:interiorbits79}.  The reference implementation rejects a
resource-limited incomplete enumeration without emitting a codebook
or payload.  Thus the theorem is a finite exact construction, not a
polynomial-time coding result.  Volumetric packing and rational nets
are standard mechanisms; the dimension-uniform operational
comparison is what transfers them to the present future-use loss.

\subsection{Randomized prefix descriptions}
The fixed-length conclusion has a counterpart which permits shared
randomness and counts expected transmitted length.  This distinguishes
an index-cardinality converse from a converse for a randomized
communication protocol.

\begin{definition}[One-way prefix description]
\label{def:prefixdescription79}
The encoder is supplied $E\in\mathfrak I_d$.  A public random
variable $R$, independent of $E$, is available to both parties.
Conditionally on each public seed, the possible finite binary
messages $W$ form a prefix-free set.  Private randomization is
allowed, and the decoder returns a legal effect $C$ from $(R,W)$
and its independent private randomness.  No timing, stopping signal,
nonpublic header, or other side channel is uncharged.  The cost is
$\sup_E\mathbb E_E|W|$ and the risk is
$\sup_E\mathbb P_E\{d_N(E,C)>\delta\}$.  The public parameters
$d,N,\delta,\eta$ are not part of the message.
\end{definition}

\begin{theorem}[Randomized expected description length]
\label{thm:prefixrate79}
Let $d,N\ge1$, $0<\delta\le2^{-13}$ and $0\le\eta\le1/8$.
Write
\[
 A_d(N,\delta)=d^2\log_2(\sqrt N/\delta),\qquad
 h_2(\eta)=-\eta\log_2\eta-(1-\eta)\log_2(1-\eta).
\]
Here $0\log_2 0=0$.  The infimum $\mathcal B_\eta$ of the costs in
Definition~\ref{def:prefixdescription79} with risk at most $\eta$
satisfies
\begin{align}\label{eq:prefixrate79}
 (1-\eta)(A_d(N,\delta)-11d^2)-h_2(\eta)
 &\le\mathcal B_\eta\\
 &\le(1-\eta)(A_d(N,\delta)+d^2\log_2 25+1).\notag
\end{align}
In particular,
$\mathcal B_\eta=(1-\eta)A_d(N,\delta)+O(d^2)$,
with absolute uniform constants.  For rational $\delta$ the
achievable nonempty messages are supplied by the finite exact
interior dictionary.
\end{theorem}
\begin{proof}
Use the $4\delta$-separated packing of
Theorem~\ref{thm:interiorentropy79} and let $X$ be uniform on its
$L$ labels, independently of $R$.  A decoded effect within
$\delta$ of the true packing point identifies its label uniquely.
Choose a nearest label with a fixed tie rule otherwise.  This
produces $\widehat X$ with error probability at most $\eta$.

For completeness, if $J$ indicates a classification error, the chain
rule gives
\[
 H(X\mid\widehat X,R)
 \le H(J)+\mathbb P(J=1)\log_2(L-1)
 \le h_2(\eta)+\eta\log_2L.
\]
The final inequality uses $\eta\le1/8$ and $L\ge4$.
Since $X$ and $R$ are independent, data processing gives
\[
 (1-\eta)\log_2L-h_2(\eta)
 \le I(X;\widehat X\mid R)
 \le I(X;W\mid R)\le H(W\mid R).
\]
The decoder's private randomness is independent conditionally on
$(R,W)$ and does not increase this information.  For each public
seed, Kraft's inequality and nonnegativity of relative entropy give
$H(W\mid R=r)\le\mathbb E(|W|\mid R=r)$; integrating proves the
same conditional inequality.  These statements also apply to
countable messages of finite expected length by truncation; infinite
cost already satisfies the lower bound.  Worst-target expected
length is at least its average over the packing.  The packing
cardinality in \eqref{eq:interiorentropy79} now proves the first
inequality.  This is the usual Fano--Kraft argument, written here to
include all the stated randomization and message conventions.

For achievability, use a public coin independent of the target.
With probability $\eta$ transmit the empty word and decode to
$I_d/2$.  With the complementary probability use the deterministic
index code of Theorem~\ref{thm:interiorcodec79}, or an interior net
for real $\delta$.  For each seed the permitted set is either the
singleton empty word or the fixed-length index set, and hence is
prefix-free.  The decoder knows which case occurs from $R$, so there
is no free transmitted abort flag.  Failure can occur only on the
public abort event.  Equations~\eqref{eq:interiorcodeerror79} and
\eqref{eq:interiorcodesize79}, with integer length rounding, give
the upper bound.  The model permits the specified Bernoulli public
coin; the zero-failure finite exact dictionary does not require it.
\end{proof}

\subsection{Learning and transmitting at the joint rates}
\begin{theorem}[Jointly optimal learned interior descriptions]
\label{thm:interiorlearnedcode79}
In the ideal-control query model of
Theorem~\ref{thm:interiorlearning78}, let $d,N\ge1$ and let
$0<\delta\le2^{-13}$ be rational.  There is one common learner
on $\mathfrak I_d$ with a finite binary output and a fixed public
deterministic decoder such that
\begin{equation}\label{eq:interiorlearned79}
 \sup_{E\in\mathfrak I_d}
   \mathbb P_E\{d_N(E,C)>\delta\}\le1/8,\qquad
 M=O(d^2N\delta^{-2}),\qquad
 B=A_d(N,\delta)+O(d^2).
\end{equation}
The constants are absolute.  The acquisition uses independent
one-call input/reference pairs, followed by collective processing
of completed outputs.  No device call is used for encoding or
replay.  Any learner with these error and decoder conventions has
$M=\Omega(d^2N\delta^{-2})$ and
$B\ge A_d(N,\delta)-11d^2$.
\end{theorem}
\begin{proof}
Apply the binary estimator of \cite[Corollary~III.9]{MB67} at
operator accuracy $\epsilon=\delta/(64k)$ and failure $1/8$.
It returns a legal effect $F$ with $\|F-E\|_{\rm op}\le\epsilon$
on its good event, using $O(d^2k^2\delta^{-2})$ calls.
Let $T$ be the spectral clipping of $F$ to $[I/4,3I/4]$.
The good event and Weyl's inequality imply
$\|T-F\|_{\rm op}\le\epsilon$, whence
\[
 \|T-E\|_{\rm op}\le2\epsilon,\qquad
 d_N(E,T)\le8\sqrt N\epsilon\le\delta/8.
\]
This argument does not assume that scalar clipping is
operator-Lipschitz for a noncommuting pair.

Round the coordinates of $T$ and encode their grid point by the
finite dictionary of Theorem~\ref{thm:interiorcodec79}.  Its proof
gives $d_N(T,C)\le5\delta/16$ even though the displayed estimator
may initially be real-valued.  The resulting map from the original
classical outcome to a dictionary index is measurable: spectral
clipping is continuous, coordinate rounding is Borel with specified
ties, and the final list is finite.  Composing the final tomography
POVM with this finite partition defines a common finite-output
measurement.  It needs neither an unknown-effect-dependent
measurement nor an exact-real side channel to the decoder.
Equivalently, certified finite coordinate approximations can be used
with the rounding slack of the coding theorem when supplied by a
classical implementation.  On the good event the total operational
error is at most $7\delta/16<\delta$.  Since $k^2\le4N$, the call
bound and the payload upper bound follow.

For the call converse, a decoded word is also an estimator, so
Theorem~\ref{thm:dimensionlower78} applies.  For the length converse,
positive success probability for every target implies that the finite
range of the fixed deterministic decoder is a $\delta$-cover.
The first inequality of \eqref{eq:interiorentropy79} applies even
when those centres lie outside $\mathfrak I_d$.
\end{proof}

The last theorem is a query and transmitted-length statement.  It
inherits the ideal trusted tomography operations of \cite{MB67};
the measurable finite readout above does not assert an efficient gate
synthesis for them.  The separate finite-control theorem for our
full-body learners remains unchanged.  The interior description
bounds neither remove the endpoint logarithms on the complete body
nor optimize its $d^4$ learning upper bound.  Rather, on the same
full-dimensional interior, they identify both training and description
costs with constants uniform in dimension, horizon and accuracy.
